\documentclass[11pt]{article}
\usepackage{amsmath,amssymb,amsthm}
\newtheorem{definition}{Definition}
\newtheorem{theorem}[definition]{Theorem}
\newtheorem{proposition}[definition]{Proposition}
\title{Configuration-indexed native arena cells: an explicit enlargement}
\date{16 November 2026}
\begin{document}\maketitle
\paragraph{Audit and scope.}
Independently reviewed within the stated local-cell, relative-assembly,
feedback and observation scope; finite diagnostics are not formal verification.
Objects, full grades, primitive laws and all configurations are unchanged from
\emph{Incremental Causal Arenas and Open Control}, SHA
\texttt{b707e7144fd95498b77ac42b77efafbc903899b41d243449e6dd669822ebaa40}.
The labelled feedback interface is that of \emph{Guarded Incremental Arena
Feedback and Cost Trace}, SHA
\texttt{cddab97194d2bb15aeb7953e93ccfabbb378d5e2e4453d936a3d28ead2927710}.
Only the cell doctrine changes. In particular this does not remove the
enabled-event obstruction to a unit-preserving context realization.
The coefficient doctrines and finite/countable semantic targets remain those
of the sources. No new effectful symmetric tensor is assumed.

\section{Actual additional maps}
Write $E_P(s)$ and $O_P(s)$ for enabled events and emitted ordered ports.
Existing port inclusions along a transition are literal stable-name inclusions.
\begin{definition}
A configuration-indexed native cell $H:P\to P'$ at fixed endpoints consists of
the original configuration map $h$ preserving anchors and initials, one full
weighted-content map $\gamma$, and, for every configuration $s$,
\[
 \epsilon_s:E_P(s)\xrightarrow{\sim}E_{P'}(h s),\qquad
 \omega_s:O_P(s)\xrightarrow{\sim}O_{P'}(h s).
\]
The port bijection preserves order and labels. For every enabled event there is
an actual native row map $r_{s,e}$ with its augmented coefficient arrows.
Require the original transition and terminal-body equations, now using
$\epsilon_s(e)$ and $\omega_s$.
For $t=s[d,w]$, require
\[
 \omega_t(p)=\omega_s(p)\quad(p\in O_P(s)),\qquad
 \epsilon_t(e)=\epsilon_s(e)\quad(e\in E_P(s),\ e\ne d).       \tag{1}
\]
The surviving row functions and actual coefficient arrows at $(s,e)$ and
$(t,e)$ are equal under the literal row-persistence identifications. These
requirements apply to every raw outcome, including zero outcomes, and every
configuration, including unreachable configurations.
\end{definition}
Equation (1) is stronger than a transition square alone. It does not identify
names across configurations unrelated by port/event persistence.
No coordinate for a never-occurring identifier is added. On terminal
configurations, $\gamma$ retains its maps on all unused shapes, not only the
shape selected by that configuration.

\begin{theorem}
These cells form native hom categories $\mathsf{Arena}_{\rm loc}(A,B)$.
The original global-cell hom categories include faithfully, identically on
objects, configurations, raw maps, coefficients and grades.
\end{theorem}
\begin{proof}
For $H:P\to P'$ and $K:P'\to P''$, set
$\epsilon^{KH}_s=\epsilon^K_{h s}\epsilon^H_s$ and
$\omega^{KH}_s=\omega^K_{h s}\omega^H_s$; compose raw/coefficient arrows and
full grade maps. The transition square for $H$ says that $h(t)$ is exactly
the target transition to which (1) for $K$ applies. Thus both persistence
equations and row-arrow equalities compose. Exact probability pushforward,
body equations, identities and associativity follow on the displayed actual
data. Restrict each old global identifier function to its active/emitted
sets. This satisfies (1). Equality of restrictions determines the global
function because every identifier occurs in at least one configuration.
Thus this inclusion is faithful, not a quotient.
\end{proof}

\paragraph{Declared additional-map kernel.}
An enlarged cell comes from an old cell exactly when its local event assignments
agree wherever the same event identifier occurs and its local port assignments
agree wherever the same port identifier occurs. These conditions are
necessary and, by taking unions of the agreeing functions, sufficient.
The new doctrine stops treating cross-configuration equality of names,
without any persistence path connecting their occurrences, as mandatory
comparison data. It does not forget such names on objects, merge old cells,
erase execution history, or identify coefficient arrows.

The inclusion is not isomorphism-reflecting. Take inputs $\{0,1\}$ and one
deterministic singleton event per input, followed by the same unary terminal
output. In one arena name the events $e_0,e_1$ on the two input fibres; in
the other use the single global name $e$. Keep configurations, laws,
coefficient arrows and full grades identical. The global forgetful cell
$e_0,e_1\mapsto e$ has no global inverse. Its local inverse uses $e\mapsto
e_a$ in input fibre $a$, and is coherent along transitions. Hence it is an
isomorphism in the enlarged category. Likewise forgetting loop-deployment
identifier tags can become locally invertible when the current input
recovers the tag, provided the full grades are identical. Thus the added-map
kernel admits new isomorphisms between previously distinct global namespace
presentations, despite being injective on all old parallel cells. It does
not remove an arity mismatch in full grades or an enabled-event mismatch.

\section{Open assembly and the restored content realization}
\begin{theorem}
Original open arena composition extends to the new cells and gives a
bicategory, with the same objects, horizontal arrows, associators and units.
Pure input/output actions and pure-middle balance give the same relative
pseudopromonad construction in the enlarged cell doctrine.
The zero-event source-content realization $N\to\mathsf{Arena}_{\rm loc}$ is
locally fully faithful and preserves composition and unit by invertible
native comparisons.
\end{theorem}
\begin{proof}
At a composite configuration $z=(s,(t_p)_{p\in O_P(s)})$, a pair of cells
$H:P\to P'$ and $K:Q\to Q'$ sends
\[
 z\longmapsto
 (h s,(k t_p)_{\omega^H_s(p)}).
\]
Root events map by $\epsilon^H_s$; a child event $(p,e)$ maps to
$(\omega^H_s(p),\epsilon^K_{t_p}(e))$; an output port $(p,q)$ maps to
$(\omega^H_s(p),\omega^K_{t_p}(q))$.
All are local bijections of the required typed sets; the last preserves
lexicographic order and labels. Rows use the corresponding component arrows.

If a root event fires, all old $\omega^H_s(p)$ stay fixed by (1).
The new-port complements correspond bijectively: $\omega_t$ is a bijection
and agrees with $\omega_s$ on the old sets. Thus each new source child
initializes precisely the corresponding new target child, and $K$ preserves
its initial configuration. If a child fires only its component changes.
This proves the transition square in both cases. For a surviving event or
port, apply (1) in its component; for an unrelated component all data are
unchanged. Actual row persistence is checked in these same cases.
The full content arrow is the original weighted substitution of $\gamma_H$
and $\gamma_K$, with its full tensor coefficient arrows.

Vertical composition of these formulas agrees exactly with horizontal
composition of the pointwise composites. This proves interchange, not just
object-level assembly. Rebracketing sends the same configurations and
addressed rows to the same three-level family. At any fixed configuration
each pentagon/triangle path is the same local address function and actual
coefficient reassociation. Naturality follows from the displayed formulas.
The old global associators/unitors therefore remain valid, and coherent,
in the larger category.

Pure input pullback restricts the local maps to pulled-back configurations;
pure output action changes labels, not these maps. The source's pure-middle
balance removes/restores the same redundant interface label. Its component
formulas commute with the local copy-indexing formula above. Identity and
composite balance laws are equalities on configurations and local maps.
Thus the same relative pseudocoend descent applies, with actual new cells.

For a source content arrow $(C,f:A\to F_C B)$ take one terminal initial
configuration per input, no events, rank-labelled ports, and full grade $C$.
A cell between two such objects has its configuration map forced by anchors;
its ordered port bijection is uniquely forced; its remaining data are exactly
the source full grade/body cell. This proves local full faithfulness.
The composite of two realizations has one configuration per input. Its ports
are pairs $(i,j)$, and the comparison with the realized source composite maps
\[
 (i,j)\longmapsto \sum_{k<i} n_k+j,                         \tag{2}
\]
where $n_k$ is the arity of the $k$th inner output at that configuration.
This is an ordered bijection; its inverse unflattens using the same arities.
There are no transitions, so (1) adds no constraint across distinct inputs.
The full grade is the original substituted grade, with its canonical
coefficient comparison. Threefold flattening is the same leaf-rank map in
either order. Unit comparisons delete singleton position coordinates.
This proves naturality, pentagon and triangle on actual maps and full grades.
These comparisons also commute with the relative pure interfaces. Input
reindexing along any $u:A'\to A$ carries the entire arity profile at $a'$
to the profile at $u(a')$, so (2) is unchanged, even when $u$ is noninjective.
Final output relabelling changes no arity or position. At the middle balance
for $v:B\to B'$, the inner row chosen at outer label $b_i$ is the same row
at $v(b_i)$ on both representatives; hence the whole ordered inner-arity
profile, local flattening and full coefficient map coincide. The source
and arena unit maps use the same singleton position. Thus this is a strong
relative comparison for $N$, not only a bicategorical composition comparison.
The old obstruction, $(2,1)$ becoming rank two or rank three at different
inputs, is precisely a new map supplied by (2).
\end{proof}
This statement retains arbitrary pure interfaces and all source content cells.
It does not assert full faithfulness of the evaluated operational composite
into its source semantics for arbitrary arenas.

\section{Feedback on the enlarged doctrine}
Keep the guard, all finite compatible invocation configurations, and the full
$R_U C$ of the labelled compiler. Keep invocation edges $(p,u)$ and the
explicitly labelled endofunctor
$T_P(X)=P|_U;\operatorname{Case}^{\rm lab}_U(X)$.
\begin{theorem}
Feedback and $T_P$ are functorial on configuration-indexed cells.
The labelled native initial-algebra theorem holds for all
$T_P$-algebras in the enlarged hom category. Labelled nesting is natural
for the new cells and retains its original coherent operational isomorphism.
\end{theorem}
\begin{proof}
At a finite feedback configuration tree, map each node configuration by $h$.
Map an outgoing edge $(p,u)$ at node $v$ to
$(\omega_{s_v}(p),u_*u)$. This recursively determines a bijection of the
present nodes and their typed paths. Map each enabled local event at $v$
by $\epsilon_{s_v}$ and each exposed local output by $\omega_{s_v}$.
These define local event/output bijections for that whole configuration.
Their raw arrows are the primitive component arrows and the full grade map
is $R_u\gamma$.

When an ancestor transitions, all existing child-edge images are fixed by
(1), and every surviving event map is fixed. New children correspond by
the new-port-complement argument; their initials correspond by preservation
of initials and the guard. A descendant transition leaves the other nodes
unchanged. These observations prove the transition square and both equations
(1) for the compiled cell, including actual row-arrow persistence. They also
prove its identity and composition equations. The case construction uses
the same local maps with the supplied loop-label tags, so $T_P$ is a functor.

Let $\alpha:T_P(X)\to X$ now be any enlarged native algebra. Define the
configuration fold recursively on finite trees:
\[
 h(s,(\tau_p))=\alpha_{\rm cfg}(s,(h(\tau_p))).
\]
At that particular configuration, define a root-event map using $\alpha$,
and a descendant-event map by its child fold followed by the corresponding
local map of $\alpha$. Define port maps in the same way; compose actual row
and coefficient arrows along the finite ancestor path. This yields
bijections and exact row maps. Guarded initials give
$h(i_F(u))=i_X(u)$.
For a root transition the existing folded children are unchanged and newly
born children fold to the required initials. For a child transition use
the inductive square for that child, followed by the square for $\alpha$.
Persistence of a surviving ancestor/descendant event or port follows by
this same case analysis and (1) for $\alpha$ and the child fold. Induction
is on the old finite ancestor structure and remains valid when the next
configuration has new children.

Independently, fold every formal grade tree with the full grade map of
$\alpha$. Tensor the child-fold coefficient arrows with the root identity
and then apply its actual content arrow. This defines the map on unused
trees too. The terminal body equation follows recursively from that for
$\alpha$. Consequently the fold is a new-doctrine native algebra map.

The algebra equation forces the configuration recursion, the root-event
map, and each descendant map through its child. It therefore forces every
local event/port map and every raw/coefficient arrow, by induction on finite
trees and event depth. The full grade equation forces the same recursion
at every formal grade tree. Hence the fold is unique in the enlarged
category. This is not merely the old initiality theorem applied to its old
smaller algebra class.

For nesting, flatten a tree of coloured trees as in the source. At each
configuration the new cell maps an original node using precisely its
original local configuration. Therefore flattening before or after this
map gives the same node, local event and port functions. The full grade
comparison and coefficients are unchanged from the source. This proves
naturality for new cells and the existing higher nesting coherence.
\end{proof}
The labelled namespace has not been removed merely because the larger
doctrine sometimes admits additional isomorphisms. No policy descent
through arbitrary algebra folds is asserted.

\section{Policies and independent evaluation}
A history map uses $\epsilon_s$ at its current configuration to transport
the selected scope and uses the primitive raw functions on its outcomes.
For simultaneously enabled events, (1) and row persistence ensure that
splitting a scope in any order gives the same target history modulo the
same exchange. History maps compose exactly, including off-policy and
zero histories and histories beginning at any retained configuration.

\begin{proposition}
The original compatible-control relations, terminal evaluation, and canonical
backward policy lifting extend to the new doctrine. Canonical lifting is an
actual contravariant functor, not equality of the larger compatible-control
relations.
\end{proposition}
\begin{proof}
At a source history pull back the target scope through its current event
bijection. For the resulting tuple of row maps $r_i$, target coupling $q'$,
and source and target primitive laws $\mu_i,\mu'_i$, put
\[
 L_rq'(w_1,\ldots,w_m)=
 q'(r_1w_1,\ldots,r_mw_m)
 \prod_i\frac{\mu_i(w_i)}{\mu'_i(r_iw_i)}.                  \tag{3}
\]
Set the whole value to zero if any denominator is zero.
Exact row pushforward means a zero target fibre has only zero source mass.
Finite fibre summation proves that (3) has source marginals and pushes to
$q'$. For composable row maps, the factors telescope on positive
denominators. All other tuples have zero mass on both sides, proving
$L_rL_t=L_{tr}$; $L_{\rm id}=1$. Current-scope pullback and history transport
compose exactly by the local cell formulas. Thus these identities hold at
every control-table entry, without dividing by history probability.
All actual context arrows and full grades are retained by the selected
cell; (3) only specifies its conditional raw law.

For the larger relation, ordinary conditional pushforward still gives
$R_KR_H\subseteq R_{KH}$ and the category of compatible selected objects.
No equality of these relations, or horizontal closure of pairs of
standalone policies, follows. Prefix pushforward follows by induction,
using the actual local row maps at each history. Enabled bijections reflect
terminality, so restriction to first-terminal strata and nonnegative
countable summation give exact terminating pushforward, with tensor
coefficient arrows and full terminal body square.

Independent policies are preserved by (3): lifting a product gives the
source product. Trace exchanges are respected by (1) and actual row
persistence, so the source independent evaluator is functorial on the new
maps. Its composition comparison remains the same bijection of complete
traces, with the same coefficient tensor and full grade. Naturality follows
by the componentwise copy/event formulas already proved. Thus its strong
relative comparison survives this enlargement.
\end{proof}

\section{Assembly with well-founded silent closure and path observation}
Use the independently reviewed silent-closure source
\emph{Well Founded Silent Closure and Native Feedback}, SHA
\texttt{5889a2e4ac7eeba5409bf5b8051bb470490c6b8c8c4bfc8f78cf2280026c81d6}.
Its object admission is well-foundedness of the initial-port relation
$u\leadsto u'$ on the whole loop interface. Finite initial branching then
makes every initial invocation closure $I_P(x)$ finite. This condition,
its equivalence to finite closure at every loop input, and the exact nesting
admission law are properties of objects and unchanged by the cell enlargement.

\begin{theorem}
For admitted objects, the well-founded-silent compiler is functorial on
configuration-indexed cells. Its initial-algebra characterization holds
among all enlarged native algebras. Its nesting, observed prefix and
event/cost-layer comparisons have the same local-cell naturality, with
conditional-compatible controls where controls are required.
\end{theorem}
\begin{proof}
At the initial closure $I_P(x)$, apply $h$ to the root initial
configuration and its local port bijection $\omega_{i_P(x)}$ to the child
slots. At each such child use the same prescription at its initial
configuration. Finite recursion gives a bijection to $I_{P'}(x')$:
initials and labels are preserved, and at every matched pair of nodes
the port bijection supplies exactly all target children. This is an
isomorphism of the finite invocation indexing trees, not a claim that
the primitive configuration or raw-row maps are invertible.
The local enabled/visible-port maps on the closure are precisely those
obtained by the recursive feedback formula. Old-edge persistence (1)
ensures that an ancestor transition does not change any existing copy's
address image. A newly born loop port attaches its entire matched finite
initial closure. No event in that closure is fired or sampled. Therefore
the transition, surviving-event, old-port and actual row-arrow persistence
proofs apply, including at every unreachable configuration.

For an arbitrary local algebra $\alpha:T_P(X)\to X$, retain the finite-tree
fold already constructed. Its new initial equation is
\[
 \operatorname{fold}_\alpha(I_P(u))=i_X(u).                 \tag{4}
\]
Prove (4) by induction on the finite initial closure. Each initial child
folds to its corresponding $i_X(u_p)$. The root tuple is consequently
the initial configuration of $T_P(X)$, and $\alpha$ sends it to $i_X(u)$.
This also verifies all the newly initialized children in a paid
transition. The earlier ancestor induction supplies the local maps,
transition/persistence equations and actual coefficients; uniqueness is
still forced on every finite configuration and every full formal grade
tree by the algebra equation. Thus the proof covers the additional local
algebras rather than merely restricting the old theorem.

For nesting, finite silent closure commutes with cutting and flattening
the coloured invocation tree: both keep the same original initial
configurations at each node and the same entire finite child closures.
Local maps apply $\omega_{s_v}$ at that node before or after flattening,
so the two routes coincide. The exact admission law remains the source's
well-foundedness of the inner relation and the induced outer relation,
including paths through inner labels. No altered map doctrine weakens
that object condition.

For compatible complete controllers, map a selected scope through the
current local event bijection and its raw tuple map. Induction on the
number of selected steps gives strict prefix pushforward. The transition
map commutes with truncation, and enabled bijections reflect terminality,
including initially terminal $T_0$. Absorbed terminal continuations are
unique and carry unchanged finite coefficients. Thus observed history/
configuration decoders and partial terminal bodies commute, with the
same full grade and actual augmentation deletion arrows.
Every paid event is mapped to one paid event; the finite initial closure
performs none. Hence complete histories preserve event count, and
restriction to each complete cost layer gives its native map. Structural
exploration bounds $N(a,n)$ from the silent-closure source need not be
equal under arbitrary native maps. The layer comparison is defined on
complete raw traces, each fully captured by its own admitted bound, not
by identifying distinct finite horizons. This is the same qualification
needed for the source's event/macro cost comparison.
\end{proof}

More explicitly, for \emph{Native Prefix Paths and Divergence}, SHA
\texttt{c62dc99887cb402b46b35cc15876a06891fa2788906fba89579a8d00e6cadbd6},
the maps just constructed satisfy its three required equations:
\[
 r'_n f_{n+1}=f_n r_n,\qquad (f_n)_*m_n=m'_n,\qquad
 f_n^{-1}(T'_n)=T_n.                                     \tag{5}
\]
The first follows by recursive history transport, the second by finite
conditional fibre summation, and the third by the current enabled-event
bijection and absorption. Hence its inverse-limit construction gives the
continuous all-raw path map, scalar probability pushforward and exact
divergence-preimage equation for these local cells. With transported
controls an operational isomorphism gives a measure-preserving
homeomorphism. Equation (1) for ports gives a compatible order isomorphism
of each corresponding path's persistent ordered port unions.
This does not supply an infinite coefficient tensor, infinite full
polynomial body, ordinary sequential stream, or forward descent of an
arbitrary controller. Actual finite prefix coefficients, full terminal
grades and complete controller tables remain separate retained data.

\paragraph{Directed loop laws in the assembled doctrine.}
Use \emph{Pure Loop Sliding and Positive Grade Obstruction}, SHA
\texttt{6b7cef8ae3af9b59c8c1892c51f814c6111d78e7314247dc08af8308bb93e8ea}.
For a fixed pure $g:V\to U$, sliding deletes the same redundant input tags
and applies $g$ to every internal label. Given a new fixed-endpoint cell
$H:P\to P'$, at an original node $s$ the two routes apply $h(s)$ and
send its port $p$ to $\omega_s(p)$, then relabel its loop label by $g$.
The recursive formulas therefore commute on all configurations and local
event/port maps. The full-grade square is the original naturality of
$R_g$ with $R\gamma$, on every unused tree and its actual coefficient
arrow. Finite silent-initial-closure recursion proves the same initial
equation on the source's exact common well-founded domain. Identity and
composite sliding are unchanged pointwise equations. Control transport,
prefixes and event counts follow from the same operational histories.

For the reviewed post-tightening comparison
\[
 \operatorname{Fb}_U(P);Q\longrightarrow
 \operatorname{Fb}_U(P;L_Q),
\]
a pair of new cells on $P,Q$ acts at each original node by its current
local maps. The configuration flattening removes only the fixed
$U$-identity copy. An edge $(p,u)$ corresponds to
$((p,*_U),u)$; either route sends it to
$((\omega_s(p),*_U),u)$ at fixed loop interface. At a $Q$ copy the two
routes apply its same local cell. This proves actual naturality of the
operational comparison. Its full-grade square is the original naturality
of $\phi_{C,D,U}$ in the whole weighted $C,D$, so actual coefficient
coherence, its composition law and directed unit law persist.

This tightening comparison also extends to well-founded silent
initialization: $L_Q$ sends its $U$ branch through the event-free identity,
whereas all $Q$ outputs are in the external interface, including at
unreachable configurations. Thus its initial loop relation is precisely
that of $P$. In each finite initial closure both sides have the same
original initial $P$ nodes, the same initial $Q$ copies at every external
port, and only the removable identity copies at loop ports. Finite
recursion proves initialization compatibility; newly born closures use
the same argument. No new paid events or initial probability draws are
inserted. Full unused-grade obstructions still prevent a reverse native
comparison in general; the local identifier enlargement does not change
them.


\section{Boundary not repaired}
The two-event/one-event and sampled-unit/zero-event obstructions remain:
local enabled bijections still reflect terminality after each transition.
More generally any functor from native normalized context rows to this
arena hom category that sends the context unit to an operational zero-event
unit and admits the actual native evaluator as a retraction is obstructed
by the native fair-row collapse to the unit. Its image must map to a
zero-enabled initial state, hence must itself be initially terminal;
its evaluated raw row is the singleton empty trace, not natively
isomorphic to a two-point fair row. Configuration-indexed names do not
change this argument.

The repair is therefore substantial but bounded: strong source content,
open assembly, labelled native feedback and backward controls coexist in
an explicitly larger cell doctrine; a strong full source interaction
embedding is not claimed. The construction does not implement the still
stronger, scope-changing contractions needed to address that obstruction.
\end{document}

